```latex
% --- Dossier 1, Question 1 ---
Réponse: On a
[cite_start]$1000\le N$, $n\le9999$, $10\le N-n\le99$ [cite: 12, 13]
[cite_start]Il reste à faire varier $N=9999,9998...$ tout en respectant ces deux contraintes : [cite: 14]
$[\begin{matrix}9999-9989=10\\ \vdots\\ 9999-9900=99\end{matrix}]$
[cite_start]$[\begin{matrix}9998-9988=10\\ \vdots\\ 9998-9899=99\end{matrix}]$ [cite: 15]
[cite_start]$[\begin{matrix}N-n=10\\ \vdots\\ N-n=99\end{matrix}]$ [cite: 16]
(90 lignes entre crochets) [cite_start][cite: 17]
(90 lignes entre crochets) [cite_start][cite: 18]
(90 lignes entre crochets) [cite_start][cite: 19]
[cite_start]$\begin{matrix}1099-1089=10\\ \vdots\\ 1099-1000=99\end{matrix}]$ [cite: 20]
(90 lignes entre crochets) [cite_start][cite: 21]
$[\begin{matrix}1098-1088=10\\ \vdots\\ 1098-1000=98\end{matrix}]$
[cite_start]$[\begin{matrix}1097-1088=10\\ \vdots\\ 1097-1000=97\end{matrix}]$ [cite: 22]
(89 lignes entre crochets) [cite_start][cite: 23]
(88 lignes entre crochets) [cite_start][cite: 24]
[cite_start]N-n 10 [cite: 25]
(k lignes entre crochets) [cite_start][cite: 26]
[cite_start]N-1000k+9 [cite: 27]
[cite_start]1011-1001 = 10 [cite: 28]
[cite_start]1011-1000 = 11 [cite: 29]
[cite_start]] [cite: 30]
(2 lignes entre crochets) [cite_start][cite: 31]
[cite_start][1010-1000 = 10] [cite: 32]
(1 ligne entre crochets) [cite_start][cite: 33]
[cite_start]Il suffit alors de compter ces lignes- et il y en a : [cite: 37]
$(1+2+\cdot\cdot\cdot+89)+((9999-1099)+1)\times90=\frac{89\times90}{2}+8901\times90$
[cite_start]$1+2+\cdot\cdot\cdot+n=\frac{n(n+1)}{2}$ [cite: 38]
[cite_start](Où l'on use de la formule [cite: 39]
[cite_start]gauche) [cite: 40]
[cite_start]pour obtenir le terme de [cite: 41]
[cite_start]Reste à factoriser le résultat par 90: la réponse E s'impose. [cite: 42]

% --- Dossier 1, Question 2 ---
Réponse: Où l'on se doute bien qu'il est inutile de calculer $3^{100}$ (pour in-
[cite_start]formation: $3^{100}>2^{100}>10^{30}$) à la main ! [cite: 49]
Comme toujours dans ce type de questions, l'examen des petites puissances four-
nit toutes les indications (On ne s'intéresse qu'au chiffre des unités, les autres
[cite_start]chiffres sont symbolisés par un carré.): [cite: 50]
[cite_start]$3^{0}=1$ [cite: 51]
[cite_start]$3^{1}=3$ [cite: 52]
[cite_start]$3^{2}=9$ [cite: 53]
[cite_start]$3^{3}=27=\Box7$ [cite: 54]
[cite_start]Continuons: [cite: 55]
[cite_start]$3^{4}=3\times3^{3}=3\times\Box7=\Box1$ [cite: 56]
[cite_start]$3^{5}=3\times3^{4}=3\times\Box1=\Box3$ [cite: 57]
[cite_start]$3^{6}=3\times\Box3=9$ [cite: 58]
[cite_start]$3^{7}=3\times\Box9=\Box7$ [cite: 59]
[cite_start]$3^{8}=3\times\Box7=\Box1$ [cite: 60]
[cite_start]etc. [cite: 62]
[cite_start]Mais le plus remarquable est que $3^{4}=\Box1$ [cite: 63]
[cite_start]On retrouve bien $13^{8}=3^{4}\times3^{4}=\Box1\times\Box1=\Box1$. [cite: 64]
[cite_start]De même $3^{12}=3^{4}\times3^{4}\times3^{4}=(3^{4})^{3}=\Box$ 1 [cite: 65]
[cite_start]$3^{16}=\Box1$ [cite: 66]
[cite_start]etc. [cite: 67]
Et comme $100=4\times25$, on a $3^{100}=(3^{4})^{25}=\Box1$.
[cite_start]Réponse C. [cite: 68]

% --- Dossier 1, Question 3 ---
[cite_start]Réponse: Examinons le produit dizaine par dizaine (la notation du 2. est [cite: 75]
[cite_start]conservée): [cite: 76]
[cite_start]$2\times4\times6\times8=\Box4$ [cite: 77]
[cite_start]$12\times14\times16\times18=\Box2\times\Box4\times\Box6\times\Box8=\Box4$ [cite: 78]
[cite_start]De même : [cite: 79]
[cite_start]$22\times24\times26\times28=\Box4$ [cite: 80]
[cite_start]etc. [cite: 81]
Le chiffre recherché est celui des unités de $4^{10}$ (car il y a dix dizaines !)
[cite_start]Or: [cite: 82]
[cite_start]$4^{10}=4\times(4^{3})^{3}$ [cite: 83]
[cite_start]et il vient facilement : [cite: 84]
[cite_start]$(4^{3})^{3}=64^{3}=\Box4^{3}=\Box4$ [cite: 85]
[cite_start]Finalement : [cite: 86]
[cite_start]$4^{10}=4\times\Box4=\Box6$ [cite: 87]
[cite_start]Réponse D. [cite: 88]

% --- Dossier 1, Question 4 ---
[cite_start]Réponse: Nous rappelons les premières puissances de 2: [cite: 97]
[cite_start]$2^{0}=1$ $2^{1}=2$, $2^{3}=8$ $2^{4}=16$ $2^{5}=32$ $2^{6}=64$, $2^{7}=128$ $2^{8}=256$ $2^{9}=512$ [cite: 98]
Et décomposons 625 en puissances de 2, comme suit:
[cite_start]$625=512+113$ [cite: 99]
[cite_start]$=512+64+49$ [cite: 100]
[cite_start]$=512+64+32+16+1$ [cite: 101]
[cite_start]$=2^{9}+2^{6}+2^{5}+2^{4}+1$ [cite: 102]
[cite_start]Ce qui peut être présenté sous forme de tableau : [cite: 103]
The following table:
"2
","9876543210
"
"625=
","1001110001
[cite_start]" [cite: 104]
[cite_start]Ainsi, 625 s'écrit 1001110001 en base 2. Réponse C. [cite: 105]

% --- Dossier 1, Question 5 ---
[cite_start]Réponse: [cite: 112]
[cite_start]Il s'agit de trouver le plus grand N tel que $7^{N}$ divise P. [cite: 113]
Remarquons que 7 est premier: 7 divise un produit si et seulement s'il en divise
[cite_start]un terme. [cite: 114]
Attention, ce n'est pas généralement vrai: penser à 6 qui divise 12 sans diviser ni
[cite_start]2, ni 3, ni 4! [cite: 115]
[cite_start]Comme $P=3\times5\times7\times...\times97\times99$ il est clair que P est divisible par 7: [cite: 116]
[cite_start]$P/7=3\times5\times(1)\times9\times11\times\cdot\cdot\cdot\times97\times99$ [cite: 117]
Peut-on continuer? [cite_start]Oui, car il y a $21=3\times7$: [cite: 118]
[cite_start]$P/7^{2}=3\times5\times(1)\times9\times\cdot\cdot\cdot\times19\times(3\times1)\times23\times\cdot\cdot\cdot\times97\times99$ [cite: 119]
[cite_start]Le prochain impair multiple de 7 est $5\times7=35$ D'où : [cite: 120]
[cite_start]$P/7^{3}=3\times5\times(1)\times9\times\cdot\cdot\cdot\times19\times(3\times1)\times23\times\cdot\cdot\cdot\times33\times(5\times1)\times37\cdot\cdot\cdot\times97\times99$ [cite: 121]
[cite_start]Attention, à l'étape suivante, nous avons $7\times7=7^{2}$: divisons deux fois par 7: [cite: 123]
[cite_start]etc. [cite: 124]
[cite_start]$P/7^{5}=3\times5\times(1)\times9\times\cdot\cdot\cdot\times33\times(5\times1)\times37\times\cdot\cdot\cdot\times47\times(1\times1)\times\cdot\cdot\cdot\times99$ [cite: 125]
[cite_start]Bien sûr, nous arrêtons quand 99 est dépassé. [cite: 126]
[cite_start]Il reste $9\times7$, 11 x 7, 13 x 7 ( $(15\times7=105>99)$ : [cite: 127]
[cite_start]On peut donc diviser P par 7 jusqu'à 8 fois : réponse E. [cite: 128]

% --- Dossier 1, Question 6 ---
Réponse: Il faut ici remarquer que 2 est premier: 2 divise un produit si et
[cite_start]seulement s'il en divise un terme. [cite: 138]
[cite_start]Aucun des termes du produit n'est 7 fois divisible par 2 car $2^{7}=128>100$. [cite: 139]
[cite_start]Seul un terme est (au plus) 6 fois divisible par 2, c'est $64=2^{6}$. [cite: 140]
[cite_start]Seuls deux termes sont (au plus) 5 fois divisibles par 2: $32=2^{5}$ et $3\times32=96$. [cite: 141]
[cite_start]Seuls quatre termes sont (au plus) 4 fois divisibles par 2 : [cite: 142]
[cite_start]$16=2^{4}$, $3\times16=48,$ $5\times16=80$ [cite: 143]
[cite_start]Seuls six termes sont (au plus) 3 fois divisibles par 2: [cite: 144]
[cite_start]$2^{3}=8,3\times8,5\times8,...,11\times8.$ [cite: 145]
[cite_start]Seuls treize termes sont (au plus) 2 fois divisibles par 2: [cite: 146]
[cite_start]$2^{2}=4,3\times4,5\times4,...,25\times4$ [cite: 147]
[cite_start]Seuls vingt-cinq termes sont (au plus) une fois divisibles par 2: [cite: 148]
[cite_start]$2,3\times2,5\times2...,49\times2$ [cite: 149]
[cite_start]Ainsi P est divisible au plus $1\times6+2\times5+3\times4+6\times3+13\times2+25\times1=97$ [cite: 150]
[cite_start]fois par 2. [cite: 151]
[cite_start]Réponse D. [cite: 152]

% --- Dossier 1, Question 7 ---
[cite_start]Réponse: : [cite: 163]
[cite_start]Posons: [cite: 164]
[cite_start]$S(1)=1$ [cite: 165]
$S(2)=11$
[cite_start]$S(3)=111$ [cite: 166]
[cite_start]$S(4)=1111$ [cite: 167]
$S(N)=111...1$ (l'entier écrit avec N chiffres '1')
[cite_start]$=1+10+10^{2}+\cdot\cdot\cdot+10^{N-1}=\frac{10^{N}-1}{10-1}$ [cite: 168]
Bien entendu :
[cite_start]$S(N+1)=1+10+10^{2}+\cdot\cdot\cdot+10^{N}=\frac{10^{N+1}-1}{9}$ [cite: 169]
[cite_start]$T(N)=1+11+111+\cdot\cdot\cdot+S(N)$ est la somme recherchée. [cite: 170]
[cite_start]Il est d'emblée évident que [cite: 171]
[cite_start]$T(N+1)=T(N)+S(N+1)$ (a) [cite: 172]
[cite_start]Considérons alors $T(N)$ comme suit: [cite: 173]
[cite_start]1 [cite: 174]
[cite_start]+11 [cite: 175]
[cite_start]+111 [cite: 176]
[cite_start]+1111 [cite: 177]
[cite_start]+111...11 [cite: 178]
[cite_start]= T(N) [cite: 179]
[cite_start]Écrivons des 0 à droite : [cite: 180]
[cite_start]0 [cite: 181]
[cite_start]+10 [cite: 182]
[cite_start]+110 [cite: 183]
[cite_start]+1110 [cite: 184]
[cite_start]+11110 [cite: 185]
[cite_start]+.111...110 [cite: 186]
[cite_start]= 10T(N) [cite: 187]
[cite_start]Puis remplaçons-les par des 1 $(n+1$ en tout): [cite: 188]
[cite_start]Ainsi : [cite: 189]
[cite_start]1 [cite: 190]
[cite_start]+11 [cite: 191]
[cite_start]+111 [cite: 192]
[cite_start]+1111 [cite: 193]
[cite_start]+11111 [cite: 194]
[cite_start]$\frac{+111...111}{=T(N+1)}$ [cite: 195]
[cite_start]$10T(N)+(N+1)\times1=T(N+1)$ (b) [cite: 196]
[cite_start]De (a) et (b) on conclut : [cite: 197]
[cite_start]$10T(N)+(N+1)=T(N)+S(N+1)$ [cite: 198]
[cite_start]Finalement : [cite: 199]
[cite_start]$9T(N)=S(N+1)-(N+1)$ [cite: 200]
[cite_start]$\dot{9}T(N)=\frac{10^{N+1}-1}{9}-(N+1)$ [cite: 201]
[cite_start]Réponse A. [cite: 202]

% --- Dossier 1, Question 8 ---
Réponse. [cite_start]Traduisons l'énoncé [cite: 213]
$\frac{a+b}{2}=c$
[cite_start]$c-a=1$ [cite: 214]
[cite_start]$a+b+c=60$ [cite: 215]
(1) [cite_start][cite: 216]
(2) [cite_start][cite: 217]
(3) [cite_start][cite: 218]
On a donc immédiatement: $a+b+c=a+b+\frac{a+b}{2}=\frac{3}{2}(a+b)=60$
[cite_start]D'où : [cite: 219]
[cite_start]$a+b=40$ [cite: 220]
(4) [cite_start][cite: 221]
[cite_start]De même [cite: 222]
[cite_start]$c-a=\frac{a+b}{2}-a=\frac{b-a}{2}=1$ [cite: 223]
[cite_start]D'où [cite: 224]
[cite_start]$b-a=2$ [cite: 225]
(5) [cite_start][cite: 226]
[cite_start]Sommer (4) et (5) donne: $2b=42$, soit [cite: 227]
[cite_start]$b=21$ [cite: 228]
(6) [cite_start][cite: 229]
[cite_start]Puis, avec (4) ou (5), on obtient : [cite: 230]
[cite_start]$a=19$ [cite: 231]
(7) [cite_start][cite: 232]
(1) [cite_start]ou (2) donne c immédiatement : [cite: 233]
[cite_start]$c=20$ [cite: 234]
[cite_start]Réponse B. [cite: 235]
(8) [cite_start][cite: 236]

% --- Dossier 1, Question 9 ---
R´ eponse. Traduisons l’´ enonc ´e
 a+b+c
 3
 = 150 (1)
 a−b = 40 (2)
 a+b
 2
 = 120 (3)
 Avec (2) et (3), on obtient :
 a−b+ (a+b) = 40+240 ()
 On a donc imm´ediatement : 
 a = 140 (4)
 Avec (2) ou (3), cela donne :
 b = 100 (5)
 Finalement, avec (1) :
 140+100+c
 3
 = 150
 soit :
 c
 3
 = 150−2403= 70
 finalement
 c = 210 (6)
 [cite_start]R´ eponse C [cite: 244]

% --- Dossier 1, Question 10 ---
R´ eponse. Toutes les r´ eponses propos ´ ees sont enti `eres : Comme la seule 
 d´ ecomposition non triviale de 55 est 55 = 5 × 11, D s’impose imm´ediatement 
 (NB : il est bien sp´ ecifi ´ e que a < b) [cite_start]: R´ eponse D . [cite: 246]

% --- Dossier 1, Question 11 ---
[cite_start]R´ eponse. [cite: 248]
Remarquons que le chiffre des unit´es de 255 est 5, ce qui disqua- 
 [cite_start]lifie toutes les suggestions, sauf D . [cite: 249]
Plus sophistiqu´e : 
 255 = 256−1 = 28 −1 = (24 −1)(24 +1) = 15×17
 [cite_start]R´ eponse D . [cite: 250]

% --- Dossier 1, Question 12 ---
[cite_start]R´ eponse. [cite: 252]
Ces trois entiers peuvent se noter b − 1,b,b + 1. La moyenne est
 donc
 b−1+b+b+1
 3
 =
 3b
 3
 = b = 6
 [cite_start]R´ eponse E. [cite: 253]

% --- Dossier 1, Question 13 ---
[cite_start]R´ eponse. [cite: 255]
En cinq minutes, soit 1/12
 `
 eme d’heure, l’aiguille des minutes se
 [cite_start]d´eplace de 360◦12= 30◦. [cite: 256]
Comme dans le mˆeme laps de temps, l’aiguille des heures 
 se d´ eplace ´ egalement (dans le m ˆ eme sens. . .), il est clair que l’angle recherch ´e est 
 [cite_start]< 30◦: R´ eponse A . [cite: 257]

% --- Dossier 1, Question 14 ---
[cite_start]R´ eponse. [cite: 259]
` 
 A 16 : 00, l’angle form ´e par les deux aiguilles est de 120 ◦(4h =
 [cite_start]un tiers de 12h). [cite: 260]
Comme l’aiguille des heures fait un tour complet en exactement 12 heures, sa
 [cite_start]vitesse est de 360◦12= 30◦` a l’heure. [cite: 261]
En 16 minutes, soit 16 /60`emes d’heure, elle
 avance donc de
 16×30◦
 60
 = 8◦
 L’aiguille des minutes est 12 fois plus rapide que celle des heures : Elle se d´eplace 
 [cite_start]dans le mˆ eme temps de 8 ×12◦ = 96◦. [cite: 262]
L’angle recherch´e vaut donc 
 120◦ +8◦ −96◦ = 32◦
 [cite_start]R´ eponse B . [cite: 263]

% --- Dossier 1, Question 15 ---
[cite_start]R´ eponse. [cite: 266]
[cite_start]` A 12 : 00, l’angle form ´e par les deux aiguilles est de 0 ◦. [cite: 267]
Comme
 l’aiguille des heures fait un tour complet en exactement 12 heures, sa vitesse an 
 [cite_start]gulaire est de 360◦12= 30◦par heure. [cite: 268]
En 12 minutes, soit 12/60` emes = 1/5` eme d’heure, elle avance donc de
 30◦
 5
 = 6◦
 De fac¸on analogue, on montre que l’aiguille des minutes (qui est douze fois plus
 rapide) avance de
 360◦
 5
 = 6×12 = 72◦
 Ainsi, ` a 12 : 12, θ vaut
 72◦ −6◦ = 66◦
 R´ ep´ etons ce raisonnement pour conna ˆıtre θ `a 15 : 15 : 
 ` 
 [cite_start]A 15 : 00, l’angle form ´e par les deux aiguilles est de 90 ◦(3h = un quart de 12h). [cite: 269]
[cite_start]En 15 minutes, soit un quart d’heure, l’aiguille des heures avance de 30◦4= 7,5◦. [cite: 270]
Dans le mˆeme temps, l’aiguille des minutes avance de 
 360◦
 4
 = 90◦
 Ainsi, ` a 15 : 15, θ vaut 90◦ +7,5◦ −90◦ = 7,5◦ θ a donc vari´e de 
 7,5◦ −66◦ = −58,5◦
 [cite_start]R´ eponse E. [cite: 271]

% --- Dossier 1, Question 16 ---
[cite_start]R´ eponse. [cite: 273]
Par d´ efinition, 10! = 1 × 2 × 3 × 4 × 5 × 6 × 7 × 8 × 9 × 10. On
 peut alors apparier tous ces facteurs :
 [cite_start]10! [cite: 274]
= (2×5×10)×(3×9)×(4×8)×(6×7)
 Ce qui donne
 10! = 27×32×42×100
 [cite_start]et ´ elimine A et B : Il faut moins deux 0 finaux. [cite: 275]
Et il n’y en a pas de troisi` eme car 10 ne divise pas 27 ×32×42 (ce n’est mˆeme 
 [cite_start]pas divisible par 5) : R´ eponse C . [cite: 276]

% --- Dossier 1, Question 17 ---
[cite_start]R´ eponse. [cite: 279]
[cite_start]La d´ ecomposition en facteurs premiers de 10 est ais ´ ee : 10 = 2×5. [cite: 280]
Ainsi,
 N = (2×5)×(2×5)2 ×(2×5)3 ×··· ×(2×5)99 ×(2×5)100
 Soit,
 N = (2×22 ×23 ×··· ×299 ×2100)×(5×52 ×53 ×··· ×599 ×5100)
 Soit,
 [cite_start]N = 21+2+···+99+100 ×51+2+···+99+100 [cite: 281]
Reste `a appliquer nos formules : 
 1+2+···+99+100 =100×1012= 5050
 Et de conclure :
 N = 25050 ×55050
 [cite_start]R´ eponse E. [cite: 282]

% --- Dossier 1, Question 18 ---
[cite_start]R´ eponse. [cite: 283]
En base sept, seuls les sept chiffres 0,1,2,3,4,5,6 sont utilis´es : 
 [cite_start]les suggestions B,C,D,E, sont ´ ecart ´ ees d’office. [cite: 284]
[cite_start]R ´ eponse A . [cite: 285]

% --- Dossier 1, Question 19 ---
[cite_start]R´ eponse. [cite: 287]
[cite_start]Raisonnons en base douze : Ici douze se note 10 ; [cite: 288]
ajouter un z´ero 
 [cite_start]`a droite marque la multiplication par douze. [cite: 289]
Or : 
 1234 = 1230+4
 [cite_start]1230 est multiple de 12, donc de 6 : Il existe q tel que 1230 = 6q. [cite: 290]
D’o`u : 
 1234 = 6q+4
 (Plus explicitement :
 1234 = 1230+4 = 123×10+4 = 123×(2×6) +4 = 6×(2×123) +4 )
 [cite_start]R´ eponse D . [cite: 291]

% --- Dossier 1, Question 20 ---
[cite_start]R´ eponse. [cite: 294]
Il faut tenir compte de l’ordre qui lie a,b, c :
 |a−b| = b−a,|b−c| [cite_start]= c−b,|a−c| [cite: 295]
= c−a
 Reste ` a faire la moyenne de ces trois termes : 
 10
 3
 =
 b−a+c−b+c−a
 3
 =
 2c−2a
 3
 D’o`u : 
 c−a = 5 , soit a = 5
 Enfin,
 b
 a
 =
 b
 5
 =
 1
 5, soit b = 1
 Comme b doit ˆ etre sup ´ erieur ` aa, il n’y a pas de solution et c’est la r´ eponse D qui
 [cite_start]est la bonne. [cite: 296]

% --- Dossier 2, Question 1 ---
R´ eponse :
 O`u l’on note x la distance que Sophie compte parcourir avec la voi- 
 [cite_start]ture. [cite: 302]
La facturation hebdomadaire du premier loueur est la suivante :
 7×50+0,1x = 350+x10
 - Supposons que Sophie compte effectuer au plus 1500 km hebdomadaires : dans
 ce cas, s’engager aupr` es du second a un co ˆut de 
 510−x10
 On a donc (pour Sophie, (C) est moins cher) :
 510−x10< 350+x10
 Soit :
 510−350 <x10+x10⇔ 160 <x5
 Finalement :
 x > 800
 -Supposons que Sophie compte effectuer plus de 1500 km : dans ce cas, (C) fac 
 turera 510 euros, tout en restant moins cher que (L) :
 [cite_start]510 < 350+x10⇔ 160 <x10⇔ 1600 < x [cite: 303]
Conclusion : on a :
 800 < x 6 1500 ou bien x > 1600
 Reste ` a diviser par 7 pour faire la moyenne journali `ere (Attention ! Ne pas cocher 
 B : c’est la distance rapport´ ee` a la journ ´ ee qui est demand ´ee.) : 
 800/7 6 x 6 1500/7 ou bien x/7 > 1600/7
 [cite_start]x/7 peut ˆ etre proche de 800 /7 = 114,... , ce qui ´ elimine les options C ,D . [cite: 304]
[cite_start]x n’est pas n´ ecessairement plafonn ´ ee : On rejette E. Seule la r´ eponse A convient. [cite: 305]

% --- Dossier 2, Question 2 ---
R´ eponse : Soient N(resp. n) le nombre de parties gagn´ees par Bernard (resp. 
 [cite_start]Alice). [cite: 309]
Il est clair que : N +n = 10 et que le gain de Bernard est ´ egal `a son ’chiffre 
 d’affaire’ moins les pertes i.e N −n = 2. Ce qui fait un syst`eme : 
  N +n = 10N −n = 2
 Sommer les deux lignes donne : 2N = 12 ⇒ N = 6.
 Bernard a gagn´e 6 fois sur 10. 
 [cite_start]R´ eponse A . [cite: 310]

% --- Dossier 2, Question 3 ---
R´ eponse : Pendant les dix premi` eres minutes (un sixi `eme d’heure), le pre- 
 mier train parcourt 120/6=20 km : il se retrouve alors en A’ au moment o`u son 
 [cite_start]homologue part de B (AA’=20 km, A’B=120 km). [cite: 316]
Les trains se recontrent T h
 plus tard, en R.
 Pour le premier train, on a :
 A0R
 T
 = 120 ⇔ A0R = 120T
 Pour le second :
 [cite_start]O` u l’on utilise la relation BR = A0B−A0R = 120−A0R. [cite: 317]
BR
 T
 = 80 ⇔ 120−A0R = 80T ⇔ A0R = 120−80T
 Par suite :
 120T = 120−80T
 et l’on obtient :
 T =
 120
 200
 =
 6
 10
 [cite_start]Reste ` a ajouter le sixi `eme d’heure : 610+16=4660. [cite: 318]
Le compte s’est fait en heures : en minutes, cela donne 92120×60 = 46.
 [cite_start]La rencontre a lieu ` a 20 h33 : r´ eponse . [cite: 319]
[cite_start]D . [cite: 320]

% --- Dossier 2, Question 4 ---
R´ eponse : Pendant les quinze premi`eres minutes (un quart d’heure), le pre- 
 mier train parcourt 120/4=30 km : il se retrouve alors en A’ (AA’=30 km, A’B=120
 [cite_start]km). [cite: 324]
[cite_start]Les trains se recontrent T h plus tard, en R : Soit D=A’R. [cite: 325]
Reste `a former un 
 [cite_start]syst` eme de deux ´equations. [cite: 326]
Pour le premier train, on a : 
 D
 T
 = 120 ⇔ T =
 D
 120
 Pour le second :
 120−D
 T
 = 80 ⇔ T =12080−D80i.e :
 Par suite :
 D
 120
 =
 120
 80
 −
 D
 80⇔
 D
 120+
 D
 80
 =
 120
 80
 80D+120D = 1202 = 14 400 (NB : 122 = 144)
 et l’on obtient :
 D =
 14 4 60 60
 2 60 60
 = 72
 [cite_start]Reste ` a ajouter les trente premiers kilom ` etres : la valeur recherch ´ee est 102 km ; [cite: 327]
[cite_start]et la r´ eponse est D . [cite: 328]

% --- Dossier 2, Question 5 ---
R´ eponse : Dire que la distance entre les deux voitures peut ´egaler SP=130 
 m signifie que si l’une d’elles part de S, la seconde part de P. Et ce ` a la m ˆeme 
 vitesse (celle du cˆ able) : Les deux voitures se croisent ` a mi-parcours, donc `a “mi- 
 [cite_start]d´ enivel ´ e”. [cite: 332]
L’inclinaison ´etant de 30 ◦ par rapport ` a l’horizontale, le d ´ enivel ´e est de 
 [cite_start]130 sin 30◦ = 130/2 = 65 m. [cite: 333]
Ainsi, la rencontre a lieu 65/2 = 32,5 m en contre 
 [cite_start]bas : R´ eponse E. [cite: 334]

% --- Dossier 2, Question 6 ---
R´ eponse : Dire que la distance entre les deux voitures peut ´egaler SP=130 
 m signifie que si l’une d’elles part de S, la seconde part de P. Et ce ` a la m ˆeme 
 vitesse (celle du cˆ able) : Les deux voitures se croisent ` a mi-parcours, donc `a “mi- 
 [cite_start]d´ enivel ´e”. [cite: 339]
La pente de la colline est de : 
 100
 √3% =
 1
 √3
 =
 1/2
 √3/2
 =
 sin 30◦
 cos 30◦
 = tan 30◦
 Autrement dit, l’inclinaison est de 30◦ par rapport ` a l’horizontale : le d ´ enivel ´e
 [cite_start]est de 130 sin 30◦ = 130/2 = 65 m. [cite: 340]
Ainsi, la rencontre a lieu 65/2 = 32,5 m en
 [cite_start]contrebas : R´ eponse E. [cite: 341]

% --- Dossier 2, Question 7 ---
R´ eponse : Dire que la distance entre les deux voitures peut ´egaler SP signifie 
 que si l’une d’elles part de S, la seconde part de P. Et ce ` a la m ˆeme vitesse (celle 
 du cˆ able) : Les deux voitures se croisent ` a mi-parcours- apr ` es avoir avanc ´e sur 
 [cite_start]d=111m, donc. [cite: 347]
On a alors :
 v = 1 =d
 t
 =
 111
 t
 ⇒ t =
 111
 1
 = 111
 [cite_start]La rencontre se fait 1 min 51 s apr` es les d ´ eparts : R ´ eponse B . [cite: 348]

% --- Dossier 2, Question 8 ---
R´ eponse : O` u il est ´evidemment plus avantageux de commander le papier 
 [cite_start]par centaines de rames que par lots de vingt-cinq. [cite: 355]
. .
 Comme les boˆıtes de crayons se commandent par 12, il faut commander au mi 
 [cite_start]nimum 7×12 = 84 boˆıtes idoines. [cite: 356]
Ce qui coˆ ute 7 ×25 = 175 euros. Reste donc
 [cite_start]1000−175 = 825 euros (au plus) `a consacrer pour l’achat de papier. [cite: 357]
[cite_start]Avec une telle somme, on peut commander 2×100 = 200 rames pour 600 euros. [cite: 358]
Reste donc 825 − 600 = 225 euros ` a d ´epenser, ce que l’on fait en commandant 
 [cite_start]2 × 25 = 50 rames -et il reste alors 25 euros dans l’enveloppe. [cite: 359]
. .que l’on utilise
 [cite_start]pour acheter un huiti`eme lot de crayons ! [cite: 360]
[cite_start]R´ eponse E. [cite: 361]

% --- Dossier 2, Question 9 ---
R´ eponse :
 [cite_start]Il s’agit bien sˆ ur d’abaisser le co ˆ ut de la d ´epollution. [cite: 370]
R´ ecapitulons les disponibilit ´es : 
 L M Me
 A A
 B B
 C C
 (NB : Il suffit de lire le tableau en ligne puis en diagonale pour s’assurer que tous
 les d´ eplacements (A un jour, B l’autre ; B un jour, C l’autre ; C deux journ ´ees etc.) 
 sont a priori possibles.)
 Une entreprise donne son nom `a son tarif journalier : 
 [cite_start]A=1000 C /jour (au plus : A facture au m2 et se limite `a 10m2 par jour.). [cite: 371]
B=2000 C /jour.
 C=2300 C /jour.
 [cite_start]Faire venir A et C (par exemple A lundi, C Mercerdi) coˆ ute A+C= 3300 C au plus. [cite: 372]
On a clairement :
 A < B < C ⇒ A+A < A+B < B+B < B+C < C +C
 A seule est d’embl´ ee´ ecart ´ ee (elle est en sous-capacit ´e), mais A&B (quatorze m2 
 [cite_start]pour B, dix pour A) conviennent : le coˆ utA+B = 3000 est maintenant le moindre. [cite: 373]
C’est une solution possible, mais est-ce la meilleure ?
 Pour le savoir, il faut regarder la seule combinaison qui nous a ´ echapp ´e : A & C. 
 Et en commandant 6 m2 ` a A, et 18 m2 ` a C, on peut en deux journ ´ ees d ´epolluer 
 les 24 m2 pour 2300+600=2900 C.
 [cite_start]R´ eponse B . [cite: 374]

% --- Dossier 2, Question 10 ---
R´ eponse : Si une machine A coˆ ute x, une machine B coˆ ute 1 ,5x (elle est 50
 [cite_start]% plus on´ ereuse). [cite: 380]
De plus, une machine A tirant 2 pages `a la minute, une B tire 
 [cite_start]1,5×2 = 3 pages/min (car elle est 50 % plus rapide qu’une A). [cite: 381]
En notant a (resp. b) le nombre de machines A (resp. B) du futur parc, on ob 
 tient les formules :
 a+b = taille du parc
 ax+1,5xb = (a+1,5b)x = coˆut
 2 pages/min×a+3 pages /min×b = productivit´e totale 
 Le probl`eme est donc le suivant : trouver a,b tels que 
 2a+3b > 20
 avec a+1,5b minimal, puis, si plusieurs solutions se pr´ esentent, les d ´epartager en 
 [cite_start]minimisant a+b. [cite: 382]
O` u, en faisant varier a tout en minimisant b, l’on obtient :
 (1) 2×0+3×7 = 21 (coˆ ute 7 ×1,5)
 (2) [cite_start]2×1+3×6 = 20 (coˆ ute 1 ×0+6×1,5)(∗) [cite: 383]
(3) 2×2+3×6 = 22 (coˆ ute 6 ×1,5)
 (4) 2×3+3×5 = 21 (coˆ ute 5 ×1,5)
 (5) 2×4+3×4 = 20 (coˆ ute 4 ×1,5)
 (6) 2×5+3×4 = 22 (coˆ ute 4 ×1,5)
 (7) 2×6+3×3 = 21 (coˆ ute 3 ×1,5)
 (8) 2×7+3×2 = 20 (coˆ ute 1 +2×1,5)
 (9) 2×8+3×2 = 22 (coˆ ute 2 +2×1,5)
 (10) 2×9+3×1 = 21 (coˆ ute 3 +1,5)
 (11) 2×10+3×0 = 20 (coˆute 4) 
 [cite_start]Il est clair que (8) et (11) repr´ esentent un moindre co ˆut, mais c’est (8) qui [cite: 384]
mini- 
 mise la taille du parc (a+b = 9, contre a+b = 10 en (11) ) : Il suffit de commander
 une A et deux B.
 [cite_start]R´ eponse B . [cite: 385]
∗ : le calcul des coˆ uts tient compte des six imprimantes A qui sont d´ ej` ad´etenues 
 [cite_start]par la soci´ et´e ! [cite: 386]

% --- Dossier 2, Question 11 ---
[cite_start]R´ eponse. [cite: 389]
C’est E, cons´ equence du fait que la moyenne de la somme est la
 [cite_start]somme des moyennes. [cite: 390]
On peut retrouver ce r´esulat au cas par cas : 
 [cite_start]Notons S le montant total des fonds : les deux desks se le partagent ` a part ´egales. [cite: 391]
Le premier desk rapporte donc (1+2%)S2.
 [cite_start]Notons t le taux (en %) de profit du second desk : Celui-ci rapporte (1+t %)S2. [cite: 392]
Globalement, le rapport est donc de :
 (1+2%)S2+ (1+t %)S2= (1+4%)S i.e.
 (1+2%) [cite_start]+ (1+t) = 2(1+4%) i.e. [cite: 393]
2%+t%= 8%
 Finalement :t=6
 [cite_start]R´eponse E. [cite: 394]

% --- Dossier 3, Question 1 ---
[cite_start]R´ eponse. [cite: 398]
Attention au pi`ege : dire que les boules ne sont pas toutes les deux 
 rouges signifie qu’au plus l’une d’elle est rouge, et non pas qu’aucune des deux
 [cite_start]n’est rouge ! [cite: 399]
- L’id´ ee est ici de cerner d’abord l’ ´ ev´ enement compl ´ementaire 
 (i.e { avoir deux billes rouges}) de celui mentionn´ e par l’ ´ enonc ´e
 (i.e { avoir au plus une boule rouge}) car il est plus simple `a se figurer, les deux 
 ´ 
 ev
 ´ enements ´ etant li ´es par la relation : 
 P({avoir au plus une boule rouge}) = 1−P({avoir deux billes rouges}) (∗)
 - On a affaire ` a deux tirages (un par bo ˆıte) : on r´ ecup `ere finalement deux billes, 
 [cite_start]l’une venant de la [cite: 400]
premi` ere bo ˆıte, l’autre de la seconde.
 Mais ces deux ´ epreuves sont ind´ ependantes (l’issue de la premi`ere n’influence 
 pas la seconde, et r´eciproquement) 
 Math´ ematiquement, l’ind ´ ependance se traduit par l’ ´ egalit ´e : 
 P(E1,2) = P(E1)×P(E2) (∗∗)
 o
 `
 u : 
 E1 = {une bille rouge sort de la premi` ere bo ˆıte}
 E2 = {une bille rouge sort de la seconde boˆıte}
 E1,2 = {avoir finalement deux billes rouges} = E1 ∩E2
 D´ eterminons P(E1) :
 [cite_start]La premi` ere bo ˆıte contient 1+2+3+4=10 billes {b1,...,b10} (une rouge b1 et [cite: 401]
9
 billes b2,...b10, d’autres couleurs) : il y a donc 10 r´esulats possibles au premier 
 [cite_start]tirage. [cite: 402]
Mais seule la rouge nous int´ eresse. . .reste ` a appliquer la d ´efinition : 
 [cite_start]Probabilit´ e=nombre de cas favorablesnombre de cas possibles [cite: 403]
Ce qui donne :
 P(E1) = 110
 De fac¸on rigoureusement analogue, on montre que P(E2) = 410.
 Reste ` a utiliser (∗∗) :
 P(E1 ∩E2) = 110×410=4100
 Puis (∗) :
 P({avoir au plus une boule rouge}) = 1−P(E1,2) = 1−4100=96100
 [cite_start]C’est la r´ eponse B . [cite: 404]

% --- Dossier 3, Question 2 ---
[cite_start]R´ eponse. [cite: 407]
Comme toujours :
 Probabilit´ e=nombre de cas favorablesnombre de cas possibles
 Chaque main possible est une combinaison de 5 cartes prises parmi 32 : il y en a
 [cite_start]exactement C532 =32!5!27!. [cite: 408]
Concernant les mains favorables : il y a exactement huit
 [cite_start]hauteurs et autant de carr´ es. [cite: 409]
[cite_start]Une fois les quatres cartes du carr ´ e “r ´ eserv ´ees”, il [cite: 410]
[cite_start]reste 32-4= 28 cartes pour faire la cinqui` eme. [cite: 411]
On a donc 8 ×28 donnes favorables.
 La probabilit´ e recherch ´ee est donc 
 8×28
 C532
 =
 8×28×27!×5!
 [cite_start]32! [cite: 412]
=
 8×28!×5!
 32!
 [cite_start]R´ eponse D . [cite: 413]

% --- Dossier 3, Question 3 ---
R´ eponse :
 Comme toujours :
 Probabilit´ e=nombre de cas favorablesnombre de cas possibles
 Chaque main possible est une combinaison de 5 cartes prises parmi 52 : il y en
 [cite_start]a exactement C552 =52!5!47!. [cite: 415]
Concernant les mains favorables : il y a exactement
 [cite_start]13 hauteurs (A,2,3,. . .,10,V,D,R)- et autant de carr´es. [cite: 416]
Une fois les quatres cartes 
 [cite_start]du carr´ e “r ´ eserv ´ ees”, il reste 52-4= 48 cartes pour faire la cinqui `eme. [cite: 417]
On a donc 
 13×48 donnes favorables.
 La probabilit´ e recherch ´ee est donc 
 13×48
 C552
 =
 13×48×47!×5!
 [cite_start]52! [cite: 418]
=
 13×48!×5!
 52!
 [cite_start]R´ eponse D . [cite: 419]

% --- Dossier 3, Question 4 ---
[cite_start]R´ eponse. [cite: 422]
Comme toujours :
 Probabilit´ e=nombre de cas favorablesnombre de cas possibles
 Chaque main possible est une combinaison de 5 cartes prises parmi 32 : il y en a
 [cite_start]exactement C532 =32!5!27!. [cite: 423]
Concernant les mains favorables : il y a exactement huit
 [cite_start]hauteurs et autant de hauteurs possibles pour une paire. [cite: 424]
Imaginons connaˆıtre la hauteur d’une paire donn´ ee et avoir `a choisir les couleurs : 
 Il y a alors exactement C24choix- (tr`efle, cœur), (carreau, pique) etc (combinaison 
 [cite_start]de deux couleurs prises parmi quatre). [cite: 425]
Une fois fix´ees les deux cartes de la paire, 
 il reste exactement 32-2= 30 cartes parmi lesquelles sont tir´ees les trois restantes : 
 [cite_start]il y a exactement C330 fac¸ons de compl´eter la main. [cite: 426]
Finalement, on obtient 
 8×C24 ×C330
 donnes favorables.
 La probabilit´ e recherch ´ee est donc 
 8×C24 ×C330
 C532
 [cite_start]R´ eponse A . [cite: 427]

% --- Dossier 3, Question 5 ---
[cite_start]R´ eponse. [cite: 430]
Comme toujours :
 Probabilit´ e=nombre de cas favorablesnombre de cas possibles
 Chaque main possible est une combinaison de 5 cartes prises parmi 52 : il y en a
 [cite_start]exactement C552 =52!5!47!. [cite: 431]
Concernant les mains favorables : il y a exactement 13
 [cite_start]hauteurs (A,2,3,. . .10,V,D,R) et autant de hauteurs possibles pour une paire. [cite: 432]
Imaginons connaˆıtre la hauteur d’une paire donn´ ee et avoir `a choisir les couleurs : 
 Il y a alors exactement C24choix- (tr`efle, cœur), (carreau, pique) etc (combinaison 
 [cite_start]de deux couleurs prises parmi quatre). [cite: 433]
Une fois r´ eserv ´ees les deux cartes de la 
 paire, il reste exactement 52-2= 50 cartes parmi lesquelles sont tir´ees les trois 
 [cite_start]restantes : il y a exactement C350 fac¸ons de compl´eter la main. [cite: 434]
Finalement, on 
 obtient
 13×C24 ×C350
 donnes favorables.
 La probabilit´ e recherch ´ee est donc 
 13×C24 ×C350
 C552
 [cite_start]R´ eponse A . [cite: 435]

% --- Dossier 3, Question 6 ---
[cite_start]R´ eponse. [cite: 438]
Notons A=Alice ` a un carr ´ e, B=Bernard a un carr ´e.
 [cite_start]On recherche la probabilit´ e de A∩B. [cite: 439]
Une fac¸on de faire est d’utiliser la notion de
 [cite_start]probabilit´ e conditionnelle. [cite: 440]
En effet : P(A∩B) = P(A|B)P(B), et P(B) se calcule
 [cite_start]ais´ ement (voir question 3.). [cite: 441]
En consid´ erant dor ´ enavant que B est r ´ ealis ´e : 
 Alice rec¸oit 5 cartes prises parmi 52-5=47 : les mains possibles sont au nombre
 [cite_start]de C547. [cite: 442]
Imaginons alors une main favorable :
 [cite_start]Il faut choisir une hauteur parmi 12 (Bernard a d´ ej` a un carr ´ e). [cite: 443]
Reste une derni `ere 
 carte, prise parmi 47-4=43. [cite_start]On a donc 12×43 donnes favorables. [cite: 444]
Finalement, P(A|B) = 12×43C547et la probabilit´ e recherch ´ee est 13×48×12×43C552 ×C547
 [cite_start]R´ eponse A . [cite: 445]

% --- Dossier 3, Question 7 ---
[cite_start]R´ eponse. [cite: 449]
Attention, au poker, la double paire n’est pas celle du math´ematicien 
 (Quel joueur vous dirait que le carr´e d’as est une simple double paire ?) [cite_start]! [cite: 450]
Chaque main possible est une combinaison de 5 cartes prises parmi 32 : il y en a
 [cite_start]exactement C532. [cite: 451]
Aux mains favorables :
 (NB : Le jeu `a 32 cartes mobilise 4 couleurs et 8 hauteurs) 
 [cite_start]Une telle main peut se voir comme la donn´ee d’une carte et d’une double paire. [cite: 452]
Concernant celle-ci :
 On choisit simultan´ ement les hauteurs des deux paires : cela ´ equivaut `a former 
 une combinaison de 2 hauteurs (une par paire) prises parmi 8 (ou encore `a un 
 tirage sans remise de deux hauteurs prises parmi 8) : Il y a C28combinaisons pos 
 [cite_start]sibles. [cite: 453]
[cite_start]On s´ electionne alors les couleurs d’une paire donn ´ eei.e. [cite: 454]
on forme une combinai 
 son de 2 couleurs prises parmi 4 (i.e. un tirage sans remise de etc.) : il y a C24
 [cite_start]possibilit´es. [cite: 455]
Nous en sommes ` aC28 ×C24 ×C24cas favorables, et il reste ` a consid ´ erer la derni `ere 
 carte :
 Si la carte avait une hauteur que l’on retrouve dans la double paire, Alice aurait un
 [cite_start]full (trois cartes de mˆeme hauteur et une paire) : il n’y a que 6 hauteurs favorables. [cite: 456]
[cite_start]Reste la couleur : chacune des quatre couleurs est favorable. [cite: 457]
Conclusion : les mains favorables sont au nombre de
 C28 ×(C24)2 ×6×4
 La probabilit´ e recherch ´ee est donc de 
 C28 ×(C24)2 ×24
 C532
 [cite_start]R´ eponse A . [cite: 458]

% --- Dossier 3, Question 8 ---
R´ eponse. Remarquons que les propositions sont list´ ees par ordre d ´ecroissant. . [cite_start]. [cite: 461]
Il y a 50% de chances d’obtenir au moins six “pile” (et 50% de chances d’en
 [cite_start]obtenir entre z´ero et cinq). [cite: 462]
Nous pouvons donc affirmer que : 
 P(Au moins sept “pile”.) = 50% −P(Exactement six “pile”.) < 50%
 [cite_start]Ce qui ´ elimine A,B d’embl ´ee. [cite: 463]
Reste ` a´ etablir, ou du moins estimer P(Exactement six “pile”.) :
 Il y a de multiples fac¸ons de faire (exactement) six fois “pile” : six “pile” sui 
 [cite_start]vis de cinq “face” en est une. [cite: 464]
Sa probabilit´ e est de 1 /2 × ··· × 1/2 = (1/2)11 Il
 est alors clair que P(Exactement six “pile”.) > (1/2)11 ce qui abaisse la probabi 
 lit´ e recherch ´ ee` a strictement moins de 50% −(1/211) Notre pr´eambule permet de 
 [cite_start]conclure : R´ eponse E. [cite: 465]

% --- Dossier 3, Question 9 ---
[cite_start]R´ eponse. [cite: 467]
Il y a 50% de chances d’obtenir au moins six “pile” (et 50% de
 [cite_start]chances d’en obtenir entre z´ero et cinq). [cite: 468]
Nous pouvons donc affirmer que : 
 P(Au moins sept “pile”.) = 50% −P(Exactement six “pile”.)
 [cite_start]L’ensemble des tirages peut se pr´esenter sous forme de tableau ; [cite: 469]
par exemple : 
 tirage n◦ 1 2 . . . 10 11
 r
 ´ esultat pile face . [cite_start]. [cite: 470]
. pile pile
 Rappelons qu’ils sont tous ´ equiprobables. Autrement dit, ils ont tous la m ˆeme 
 [cite_start]probabilit´ e d’occurence que la s ´ erie “face”, . [cite: 471]
. ., “face”- cf. tableau suivant :
 tirage n◦ 1 2 . . [cite_start]. [cite: 472]
10 11
 r
 ´ esultat face face . . [cite_start]. [cite: 473]
face face
 [cite_start]Un tel ´ ev´ enement est de probabilit ´ e 1 /2× ··· ×1/2 = (1/2)11 ; [cite: 474]
et n’importe quel
 [cite_start]tirage d’exactement six “pile” est donc de mˆ eme probabilit ´ e 1 /2 × ··· × 1/2 = [cite: 475]
(1/2)11.
 Remarquons que la seule donn´ ee des “pile” suffit ` a d ´ efinir une s ´erie de tirages (si 
 [cite_start]ce n’est pas “pile”, c’est donc “face”)- et r´eciproquement. [cite: 476]
Combien y a t-il de fac¸on de disperser exactement six “pile” parmi les onze cases
 [cite_start]du tableau ? [cite: 477]
Et bien, C611. D’o`u : 
 P(Exactement six “pile”.) = C611(1/211)
 Et de conclure :
 P(Au moins sept “pile”.) = 50% −C611 1/211 
 [cite_start]R´ eponse B . [cite: 478]

% --- Dossier 3, Question 10 ---
[cite_start]R´ eponse. [cite: 480]
Il y a 50% de chances d’obtenir au moins six “pile” (et 50% de
 [cite_start]chances d’en obtenir entre z´ero et cinq). [cite: 481]
Nous pouvons donc affirmer que : 
 P(Au moins sept “pile”.) = 50% −P(Exactement six “pile”.)
 [cite_start]L’ensemble des tirages peut se pr´esenter sous forme de tableau ; [cite: 482]
par exemple : 
 tirage n◦ 1 2 . . . 10 11
 r
 ´ esultat pile face . [cite_start]. [cite: 483]
. pile pile
 Rappelons qu’ils sont tous ´ equiprobables. Autrement dit, ils ont tous la m ˆeme 
 [cite_start]probabilit´ e d’occurence que la s ´ erie “face”, . [cite: 484]
. ., “face”- cf. tableau suivant :
 tirage n◦ 1 2 . . [cite_start]. [cite: 485]
10 11
 r
 ´ esultat face face . . [cite_start]. [cite: 486]
face face
 [cite_start]Un tel ´ ev´ enement est de probabilit ´ e 1 /2× ··· ×1/2 = (1/2)11 ; [cite: 487]
et n’importe quel
 [cite_start]tirage d’exactement six “pile” est donc de probabilit´ e 1 /2×··· ×1/2 = (1/2)11. [cite: 488]
[cite_start]Remarquons que la seule donn´ ee des “pile” suffit ` a d ´ efinir une s ´erie de tirages (si [cite: 489]
[cite_start]ce n’est pas “pile”, c’est donc “face”)- et r´eciproquement. [cite: 490]
Combien y a t-il de fac¸on de disperser exactement six “pile” parmi les onze cases
 [cite_start]du tableau ? [cite: 491]
Et bien, C611.
 [cite_start]Remarque : cela revient `a disperser exactement cinq “face ” parmi onze cases. [cite: 492]
D’o`u : 
 C511 = C611
 D’o`u : 
 P(Exactement six “pile”) = C611(1/211) = C511(1/211)
 Et de conclure :
 P(Au moins sept “pile”.) = 50% −C511 1/211 
 [cite_start]R´ eponse C . [cite: 493]

% --- Dossier 3, Question 11 ---
R´ eponse. [cite_start]L’urne contient trois boules de chaque couleur. [cite: 496]
Remarquons tout
 d’abord qu’avoir deux bleues (vertes) n’implique pas d’en avoir exactement deux :
 [cite_start]rien n’empˆeche d’en tirer trois de suite. [cite: 497]
Ce que nous recherchons est la probabilit´ e de l’ ´ ev`enement 
 E = Eb ∪Ev
 o
 `
 u : 
 [cite_start]Eb = {avoir (au moins) deux bleues}, Ev = {avoir (au moins) deux vertes} [cite: 498]
Remarquons alors qu’une telle r´ eunion est n ´ecessairement disjointe : on ne peut 
 obtenir simultan´ ement deux bleues et deux vertes, le nombre de boules tir ´ ees ´etant 
 [cite_start]limit´ e`a trois. [cite: 499]
Ainsi, 
 P(E) = P(Eb ∪Ev) = P(Eb) +P(Ev) (1)
 Remarquons enfin qu’ {avoir (au moins) deux vertes} est de mˆ eme probabilit ´e qu’ 
 {avoir (au moins) deux bleues}
 P(Eb) +P(Ev) = 2P(Eb) (2)
 (1) et (2) donnent :
 P(B) = 2P(Eb) (3)
 Reste ` a d ´ eterminer P(Eb) en utilisant la d´efinition 
 Probabilit´ e=nombre de cas favorablesnombre de cas possibles
 – L’espace des tirages s’assimile ` a celui des combinaisons de trois ´ el´ements 
 pris parmi neuf : il y a exactement C39cas possibles, que l’on ordonne en
 types d’´ ev`enements incompatibles : 
 
 [cite_start][cite: 500]
1) Une bleue, une rouge, une verte.
 2) Deux bleues exactement.
 3) Deux rouges exactement.
 [cite_start]4) Deux vertes exactement. [cite: 501]
5) Trois bleues exactement.
 6) Trois rouges exactement.
 [cite_start]7) Trois vertes exactement. [cite: 502]
[cite_start]Remarque : Ces ´ ev` enements ne sont pas deux ` a deux ´equiprobables. [cite: 503]
– Les cas favorables sont les :
 – (2). [cite_start]Choisissons les deux bleues, parmi trois : cela laisse C23 = 3 choix. [cite: 504]
[cite_start]Reste ` a choisir la derni `ere boule, parmi les six restantes. [cite: 505]
Ainsi, il y a 
 [cite_start]exactement dix-huit fac¸ons d’obtenir seulement deux boules bleues. [cite: 506]
[cite_start]– (5), qui, ´ evidemment, ne recouvre qu’un seul ´ ev`enement. [cite: 507]
Ce qui permet d’affirmer :
 P(Eb) = 18+1C39
 Et de conclure `a l’aide de (3) : 
 P(E) = 2×1984=1942
 [cite_start]R´ eponse A . [cite: 508]

% --- Dossier 4, Question 1 ---
[cite_start]R´ eponse. [cite: 510]
Si l’arrˆete est de 2, une face vaut 4. Un cube a 6 faces : son aire est 
 [cite_start]de 24. La r´ eunion de deux cubes est d’aire 2 ×24 = 48. . [cite: 511]
.` a quoi on enl `eve l’aire 
 [cite_start]des deux faces de jonction ! [cite: 512]
[cite_start]La r´ eponse est donc D . [cite: 513]

% --- Dossier 4, Question 2 ---
[cite_start]R´ eponse. [cite: 514]
[cite_start]La circonf´ erence vaut 2 πR : la doubler, c’est doubler le rayon. [cite: 515]
[cite_start]Comme le rayon est doubl´ e, la surface du disque passe de πR2` aπ(2R)2 = 4πR2. [cite: 516]
[cite_start]Le cœfficient recherch´ e est 4 : R ´ eponse A . [cite: 517]

% --- Dossier 4, Question 3 ---
R´ eponse. La premi` ere ligne de l’ ´ enonc ´e dit simplement que le cube C est 
 [cite_start]inscrit dans la boule B (seule possibilit´e d’avoir huit points de contact) ; [cite: 521]
tout 
 comme un carr
 ´ e peut ˆetre inscrit dans un cercle, avec quatre points de contact 
 [cite_start]exactement. [cite: 522]
Donc :
 Volume(B\C) = Volume(B)−Volume(C)
 Il est clair que Volume(C) = 1. De plus, si R est le rayon de B, on a :
 Volume(B) = 43πR3
 Reste ` a trouver R.
 Comme C est inscrit dans B, Pythagore montre que :
 R =√3
 2
 Finalement :
 Volume(B\C) = 43π √32!3−1 =√32π−1
 [cite_start]R´ eponse A . [cite: 523]

% --- Dossier 4, Question 4 ---
[cite_start]R´ eponse. [cite: 527]
Soient (respectivement) H = 4/√3, R = 1 les hauteur et rayon de
 ce c
 [cite_start]ˆ one. [cite: 528]
[cite_start]De plus, la distance recherch ´ ee est h, le rayon de la section du cˆone avec [cite: 529]
[cite_start]le plan d’eau est r. [cite: 530]
Avec Thal`es, on obtient : 
 R
 H
 =
 1
 (4/√3)
 =
 r
 hd’o` u : r =h(4/√3)=√34h
 On a alors :
 Volume immerg´ e=13π×r2 ×h =13π× √34h!2×h =π16h3
 Mais comme d’apr` es l’ ´ enonc ´e
 Volume immerg´ e=3√38Volume du cˆ one =3√38×13π×12 ×4√3=12π
 il vient que :
 π
 16h3 =12π
 Soit
 h3 = 8
 Finalement :
 h = 2
 [cite_start]R´ eponse C . [cite: 531]

% --- Dossier 4, Question 5 ---
R´ eponse. NB : Comme l’aire recherch´ ee est strictement inf ´ erieure `a celle 
 de D et que cette derni` ere vaut π12 = π ' 3, les r´eponses C et D sont exclues 
 (4π−1 ' 11 > 2π−1 ' 5 > π)[cite_start]. [cite: 534]
L’hypoth´ enuse de T est un diam` etre de C : Avec Pythagore, il vient que les deux
 [cite_start]autres cˆ ot´es sont de longueur √2. [cite: 535]
Comme ils sont orthogonaux :
 Aire(T) = 12×√2×√2 = 1
 D ´etant d’aire 
 Aire(D) = π12 = π
 celle recherch´ee vaut 
 Aire(D)−Aire(T) = π−1
 [cite_start]R´ eponse E. [cite: 536]

% --- Dossier 5, Question 1 ---
[cite_start]R´ eponse : Notons r un tel taux et posons t = 1 + r. [cite: 540]
Par d´efinition du taux 
 d’inflation, ce qui valait p au d´ ebut de la d ´ ecennie valait (1+r)p = t p un an plus
 [cite_start]tard, t ×t p = t2 p deux ans plus tard,. [cite: 541]
. ., t10 p dix ans plus tard-correspondant `a
 une 
 ´erosion de 10 % : 
 t10 p = (1+10%)p = 1,1p
 Ainsi :
 t10 = 1,1
 Finalement :
 t = 1+r = 1,1
 1
 10
 [cite_start]C’est la r´ eponse C . [cite: 542]

% --- Dossier 5, Question 2 ---
[cite_start]R´ eponse : Compte tenu des suggestions, cette expression a un sens. [cite: 545]
Consid´erons 
 la suite suivante :
 x0 = 1
 x1 = 1+1x0
 x2 = 1+1x1= 1+11+11=32
 x3 = 1+1x2= 1+11+1
 1+11
 = 1+23 = 1,66...
 x4 = 1+1x3= 1+11+1
 1+1
 1+11
 [cite_start]. [cite: 546]
.
 .
 xn+1 = 1+1xn(∗)
 .
 .
 .
 Il est alors manifeste que la suite x0, x1,..., xn, xn+1,... tend vers un certain r´eel 
 x = 1+11+1
 1+11+···
 > 0. En passant (∗) `a la limite, on a donc : 
 x = 1+
 1
 x
 On peut aussi faire valoir le point de vue suivant :
 Posons x = 1+11+1
 1+11+···
 > 0. On a donc :
 1
 x
 =
 1
 1+11+11+···puis :
 1+
 1
 x
 = 1+
 1
 1+11+11+···
 = x
 Dans un second temps :
 1
 x
 [cite_start]+1 = x ⇔ 1+x = x2 ⇔ x2 [cite: 547]
−x−1 = 0
 [cite_start]Des ´equations qui, a priori, ont pour solution(s) : 1±√52. [cite: 548]
[cite_start]Comme x > 0, seule 1+√52 = φ (le nombre d’or) est acceptable. [cite: 549]
[cite_start]Remarquons que 2φ−1 =√5 : R´ eponse D . [cite: 550]

% --- Dossier 5, Question 3 ---
[cite_start]R´ eponse : Remarque : un mouton a deux cˆ ot´es donc au moins quatre pattes. [cite: 556]
[cite_start]Il y a alors au moins quarantes pattes dans le pr´e. [cite: 557]
Or, il y en a au plus quarante : il 
 [cite_start]y en a donc exactement quarante. [cite: 558]
Comptons les moutons (1,2,. . .,10) et repr´esentons lesdites pattes par des ’I’ (un 
 ’I’= une patte)- nous partons de la configuration moyenne (40 pattes/10 moutons=
 4 pattes par mouton) :
 [cite_start]1 2 3 . [cite: 559]
. . . . . 9 10
 IIII IIII IIII . . . . . [cite_start]. [cite: 560]
IIII IIII
 [cite_start]Est-il possible de r´epartir autrement, les quarantes pattes ? [cite: 561]
Par exemple, le mouton 
 [cite_start]1 pourrait-il en avoir cinq ? [cite: 562]
Et bien non car cela l’obligerait ` a ’ ´echanger’ une patte 
 avec un autre (par exemple, 10) :
 [cite_start]1 2 3 . [cite: 563]
. . . . . 9 10
 IIII.I IIII IIII . . . . . . [cite_start]IIII III. [cite: 564]
[cite_start]R´ eponse D . [cite: 565]

% --- Dossier 5, Question 4 ---
R´ eponse :
 [cite_start]Consid´ erons le cercle trigonom ´ etrique : quand x va de −π ` a 0, cos x croˆıt de −1 `a [cite: 567]
1.
 Donc 1+cos x > 0 croˆıt de 0 `a 2. 
 [cite_start]Donc 11+cos xd´ ecro ˆıt de +∞ ` a 1 /2. [cite: 568]
Donc −11+cos xcroˆıt de −∞ ` a−1/2.
 Donc exp −11+cos x croˆıt de 0 de e−1/2.
 [cite_start]R´ eponse B . [cite: 569]

% --- Dossier 5, Question 5 ---
R´ eponse : On a d’embl´ee : 
 [cite_start]E 6 min{B;D} 6 max{B;D} [cite: 570]
6 A
 Comme A 6 C, on a de surcroˆıt :
 E 6 min{B;D} 6 max{B;D} 6 A 6 C
 Reste ` a remarquer que max {C;D} = C 6 E pour conclure :
 E 6 min{B;D} 6 max{B;D} 6 A 6 C 6 E
 [cite_start]R´ eponse A . [cite: 571]

% --- Dossier 5, Question 6 ---
R´ eponse : p rec¸oit x euros, une deuxi` eme personne rec¸oit 2 x euros, une
 troisi` eme 2 × 2 = 4x = 22x euros, une quatri` eme 2 × 4x = 23x euros etc. On
 peut alors utiliser la formule donnant la somme des n premiers termes d’une suite
 g
 ´ 
 eom
 ´ etrique- dans le cas pr ´esent, il existe un raccourci : 
 1+2+4+···+2n−1 = 2n −1 (n > 1)
 D’o` u, en fixant n=6 et multipliant par x
 1x+2x+4x+···+25x = (26 −1)x
 Conclusion :
 [cite_start]126 = (26 −1)x = 63x [cite: 574]
⇔ x = 2
 [cite_start]C’est la r´ eponse A . [cite: 575]

% --- Dossier 5, Question 7 ---
R´ eponse : Comme toujours avec les suites, il faut calculer les premiers
 [cite_start]termes. [cite: 578]
Ce qui donne :
 U0 = 9 ⇒ U1 = 9/10+9 = 9,9 ⇒ U2 = 9,9/10+9 = 0,99+9 = 9,99 ⇒ ...
 Il est maintenant clair que U3 = 9,999,U4 = 9,9999 etc. et que la suite tend vers
 [cite_start]9,99999··· = 10 : R´ eponse D . [cite: 579]
Attention ` a la r ´ eponse C) : la suite n’atteint pas 10 (aucun U n n’est ´ egal `a 10), 
 [cite_start]mais elle s’en rapproche “`a l’infini”. [cite: 580]

% --- Dossier 5, Question 8 ---
[cite_start]R´ eponse : tan((−1)nπ/4) = tan(±π/4)±1. [cite: 581]
Donc
 Un =
 1
 n2 −1
 =
 1
 n+1
 ·
 1
 n−1
 C’est une fraction de pˆ oles -1 et 1.
 [cite_start]Comme : max{−1; [cite: 582]
1} −min{−1; 1} = 1−(−1) = 2, on a :
 Un =
 1
 2 1n−1−1n−(−1) 
 Reste ` a sommer les Un :
 U2 =1/21 −1/23
 +U3 =1/22 −1/24
 +U4 =1/23 −1/25
 +U5 =1/24 −1/26
 +U6 =1/25 −1/27
 [cite_start]. [cite: 583]
.
 .
 .
 .
 .
 =
 3
 4
 [cite_start]R´ eponse C . [cite: 584]
```